\section{Before Jesus Christ}

At death we shall stand before Jesus Christ. He will judge the life we have actually lived: what we loved, what we did, what we refused, and what we made of his gifts. The judgment will reach the hidden heart as surely as the public deed. No favorable account of ourselves can replace the truth; no forgotten act of fidelity will be lost. The Lord before whom we shall stand is already calling us to repentance. The time in which we can answer him is being given to us now.

The Creed identifies the Judge before we attempt to imagine the judgment. He is the Son who became man for our salvation, suffered under Pontius Pilate, was crucified, rose from the dead, and ascended to the Father. His earthly life teaches us whom we shall meet. We have heard him summon sinners, receive the penitent, rebuke hypocrisy, feed the multitude, and give himself for his enemies. We can learn to trust his judgment by attending to these acts, including the rebukes which call us to change.

St.~Cyril of Jerusalem begins his catechesis on Christ's return by placing the two comings beside one another. At the first, the Lord was wrapped in swaddling clothes; at the second, he comes in glory. At the first, he endured the Cross; at the second, angels attend him. The Savior who submitted to judgment will himself judge (\emph{Catechetical Lectures} XV.1--2). Cyril teaches his hearers to confess both comings. The humility in which Christ suffered and the glory in which he will judge belong to the Lord whom they are preparing to receive in Baptism.

Before the crucifix we see what sin costs and how God loves sinners. Christ bears our redemption through obedience unto death. He gives himself so that we may live. To meditate on his judgment here is to acknowledge that the one to whom we must give an account has shed his blood for us. Repentance answers that love. We can bring a particular betrayal before him, ask forgiveness, and accept the change which forgiveness requires.

In John's Gospel, Christ speaks of his authority to judge and his power to give life. The Father has entrusted judgment to the Son; the dead will hear his voice; those in the tombs will come forth to the resurrection of life or of judgment (Jn 5:21--29). He can raise the dead, and he can raise the sinner from spiritual death. Hearing his word now is already an encounter with the Lord whose voice will finally reach every grave. We prepare by receiving that word in faith and obeying it.

St.~Thomas explains that Christ exercises judgment also in his humanity, as the head of the human family he redeems. Among the reasons Thomas gives is his kinship with us: God judges human beings through the man Christ, who shares our nature and knows our weakness without sharing our sin (\emph{Summa theologiae} III, q.~59, a.~2). Our Judge has known hunger, weariness, sorrow, abandonment, and temptation. He understands the conditions under which we struggled to obey. A person who cannot explain his suffering adequately to anyone else can entrust it to Christ without fear of being misunderstood.

Thomas also relates Christ's judicial authority to the merit of his Passion. The one who fought for divine justice and was unjustly condemned is fittingly the Judge of all (III, q.~59, a.~3). Christ's exaltation reverses the condemnation pronounced against him. His servants can therefore await his judgment even when an earthly court has condemned them unjustly. A verdict obtained through violence or fear cannot make their fidelity a fault before him.

Confidence in Jesus includes reverence for his commands. We cannot appeal to his mercy while deliberately reserving a part of life in which he may not rule. His friendship seeks our conversion. When we recognize a sin and continue protecting it from correction, we resist the Lord whose mercy we ask. We can instead acknowledge it before him and seek the help needed to abandon it. We approach as sinners asking to be made clean.

Annual meditation gives time for this examination. We can consider how we have responded to Christ during the year, including a refusal which habit has made easy to overlook. Past sins can be brought to him in repentance; past gifts can be received with thanksgiving. Where a duty has been neglected, the examination should lead to a decision about fulfilling it. The time available for that response is the life we still have.

